% Общие компоненты единого шаблона. Настройки страницы — в mietlab.cls.
% Все размеры шрифта задаются в bp: 1 bp = 1/72 дюйма, как pt в Word/PDF.
\RequirePackage{fontspec}
% Word-образец не применяет кернинг/лигатуры к этим строкам титульника.
\defaultfontfeatures{Ligatures=NoCommon,Kerning=Off}
\IfFileExists{fonts/times.ttf}{%
  \setmainfont{times.ttf}[Path=fonts/,BoldFont=timesbd.ttf,
    ItalicFont=timesi.ttf,BoldItalicFont=timesbi.ttf]
  \setsansfont{times.ttf}[Path=fonts/,BoldFont=timesbd.ttf,
    ItalicFont=timesi.ttf,BoldItalicFont=timesbi.ttf]
  \setmonofont{times.ttf}[Path=fonts/,BoldFont=timesbd.ttf,
    ItalicFont=timesi.ttf,BoldItalicFont=timesbi.ttf]
}{%
  \IfFontExistsTF{Times New Roman}{%
    \setmainfont{Times New Roman}
    \setsansfont{Times New Roman}
    \setmonofont{Times New Roman}
  }{%
    \ClassError{mietlab}{Times New Roman is required}
      {Install Times New Roman or put times.ttf, timesbd.ttf,
       timesi.ttf, timesbi.ttf into the fonts/ folder. No substitute is used.}
  }
}
\RequirePackage{polyglossia}
\setdefaultlanguage{russian}
\setotherlanguage{english}
% Не задаём отдельные семейства для кириллицы: основной шрифт уже её содержит.
% Это сохраняет и поддержку локальной папки fonts/ в Overleaf.

\makeatletter
\renewcommand\normalsize{%
  \@setfontsize\normalsize{13bp}{19.5bp}%
  \setbox\strutbox\hbox{\vrule height13bp depth6.5bp width0pt}%
  \abovedisplayskip=0bp \belowdisplayskip=0bp
  \abovedisplayshortskip=0bp \belowdisplayshortskip=0bp}
\let\tiny\normalsize
\let\scriptsize\normalsize
\let\footnotesize\normalsize
\let\small\normalsize
\let\large\normalsize
\let\Large\normalsize
\let\LARGE\normalsize
\let\huge\normalsize
\let\Huge\normalsize
\renewcommand\@makefnmark{\hbox{\raisebox{.45ex}{\normalfont\normalsize\@thefnmark}}}
\makeatother
\normalsize
\linespread{1}
\setlength\parindent{0bp}
\setlength\parskip{0bp}
\setlength\topskip{13bp}
\setlength\emergencystretch{3em}
\raggedbottom
\frenchspacing
\clubpenalty=10000
\widowpenalty=10000

\RequirePackage{ragged2e}
\RequirePackage{array}
\RequirePackage{longtable}
\RequirePackage{graphicx}
\RequirePackage{float}
\RequirePackage{xcolor}
\RequirePackage{caption}
\RequirePackage{listings}
\RequirePackage{textcomp}
\RequirePackage[breakable,listings]{tcolorbox}
\RequirePackage{needspace}
\RequirePackage{enumitem}
\RequirePackage{hyperref}
\hypersetup{hidelinks,unicode=true,pdfcreator={XeLaTeX / mietlab}}
\urlstyle{same}
\setlength\intextsep{0bp}
\setlength\textfloatsep{0bp}
\setlength\floatsep{0bp}
\setlist{font=\normalfont\normalsize,topsep=0bp,partopsep=0bp,
  itemsep=0bp,parsep=0bp,leftmargin=1.25cm}

% Заголовки — слева, полужирные, с тем же кеглем и интервалом.
% Нет автоматически добавляемых номеров, отступов и пустых строк.
\makeatletter
\renewcommand\section{\@startsection{section}{1}{0pt}
  {0bp}{1sp}{\normalfont\normalsize\bfseries}}
\renewcommand\subsection{\@startsection{subsection}{2}{0pt}
  {0bp}{1sp}{\normalfont\normalsize\bfseries}}
\renewcommand\subsubsection{\@startsection{subsubsection}{3}{0pt}
  {0bp}{1sp}{\normalfont\normalsize\bfseries}}
\makeatother
\setcounter{secnumdepth}{0}
\newcommand\LabHeading[1]{\par\Needspace{39bp}\section*{#1}}
\newcommand\LabLabel[2]{\par\noindent\textbf{#1:} #2\par}
\newcommand\LabNewPage{\clearpage}

% L принимает полную ширину столбца между линиями сетки, включая поля.
% Для общей ширины \linewidth вычтите одну толщину линии в последнем L.
\newcolumntype{L}[1]{>{\normalfont\justifying\arraybackslash
  \hyphenpenalty=10000\exhyphenpenalty=10000
  \setlength\parindent{0pt}\setlength\parskip{0pt}}
  p{\dimexpr#1-2\tabcolsep-\arrayrulewidth\relax}}
\newcolumntype{C}[1]{>{\normalfont\centering\arraybackslash
  \hyphenpenalty=10000\exhyphenpenalty=10000}
  p{\dimexpr#1-2\tabcolsep-\arrayrulewidth\relax}}
\newcommand\LabTableHead[1]{\textbf{#1}}
\newenvironment{labtable}[1]{%
  \par\begingroup
  \normalsize
  \setlength\tabcolsep{5bp}%
  \setlength\arrayrulewidth{0.5bp}%
  \setlength\extrarowheight{0bp}%
  \renewcommand\arraystretch{1}%
  % Строка: 16 bp до базовой линии и 9.5 bp после неё, включая линии.
  \setbox\strutbox\hbox{\vrule height15.75bp depth9.25bp width0pt}%
  \setlength\LTpre{3bp}\setlength\LTpost{3bp}%
  \setlength\LTleft{0bp}\setlength\LTright{\fill}%
  \begin{longtable}{#1}%
}{%
  \end{longtable}\endgroup\par
}

% Подпись рисунка, как во втором PDF: по центру, снизу, обычный шрифт.
\DeclareCaptionFont{labtext}{\normalfont\normalsize}
\DeclareCaptionLabelSeparator{labdash}{\space—\space}
\captionsetup{font=labtext,labelfont=labtext,textfont=labtext,
  justification=centering,singlelinecheck=false,labelsep=labdash,
  skip=7.25bp,belowskip=0bp,position=bottom}
\gappto\captionsrussian{\renewcommand\figurename{Рисунок}}
% Одна пустая строка перед изображением и после подписи.
% Эти отступы находятся внутри блока рисунка и переносятся вместе с ним.
\newlength\LabFigureOuterSkip
\setlength\LabFigureOuterSkip{19.5bp}
\newcommand\LabFigure[4][0.94\linewidth]{%
  \begin{figure}[H]
    \vspace*{\LabFigureOuterSkip}%
    \centering
    \includegraphics[width=#1,keepaspectratio]{#2}\par
    \caption{#3}\label{#4}
    \vspace*{\LabFigureOuterSkip}%
  \end{figure}%
}

% Даже код — Times New Roman 13 bp, интервал 19.5 bp.
% Пропорциональный шрифт обязателен по запросу; пробелы сохраняются.
\definecolor{labcodebg}{RGB}{245,245,245}
\newlength\LabCodePadding
\setlength\LabCodePadding{4mm}
\newlength\LabCodeOuterSkip
\setlength\LabCodeOuterSkip{19.5bp}
\lstdefinestyle{labcode}{%
  basicstyle=\normalfont\normalsize,
  keywordstyle=\normalfont\normalsize,
  commentstyle=\normalfont\normalsize,
  stringstyle=\normalfont\normalsize,
  identifierstyle=\normalfont\normalsize,
  numberstyle=\normalfont\normalsize,
  % Фон и внутренние поля задаёт контейнер labcode, а не сам листинг.
  backgroundcolor={},
  columns=fullflexible,keepspaces=true,showstringspaces=false,upquote=true,
  breaklines=true,breakatwhitespace=false,breakindent=0bp,
  numbers=none,frame=none,tabsize=4,
  aboveskip=0bp,belowskip=0bp,
  xleftmargin=0bp,xrightmargin=0bp,
  framesep=0bp,escapeinside={(*@}{@*)},
  % listings увеличивает стандартный знак *; выводим его напрямую без масштаба.
  literate={*}{{\char"002A}}1 {-}{{\char"002D}}1
    {а}{{а}}1 {б}{{б}}1 {в}{{в}}1 {г}{{г}}1 {д}{{д}}1
    {е}{{е}}1 {ё}{{ё}}1 {ж}{{ж}}1 {з}{{з}}1 {и}{{и}}1 {й}{{й}}1
    {к}{{к}}1 {л}{{л}}1 {м}{{м}}1 {н}{{н}}1 {о}{{о}}1 {п}{{п}}1
    {р}{{р}}1 {с}{{с}}1 {т}{{т}}1 {у}{{у}}1 {ф}{{ф}}1 {х}{{х}}1
    {ц}{{ц}}1 {ч}{{ч}}1 {ш}{{ш}}1 {щ}{{щ}}1 {ъ}{{ъ}}1 {ы}{{ы}}1
    {ь}{{ь}}1 {э}{{э}}1 {ю}{{ю}}1 {я}{{я}}1
    {А}{{А}}1 {Б}{{Б}}1 {В}{{В}}1 {Г}{{Г}}1 {Д}{{Д}}1 {Е}{{Е}}1
    {Ё}{{Ё}}1 {Ж}{{Ж}}1 {З}{{З}}1 {И}{{И}}1 {Й}{{Й}}1 {К}{{К}}1
    {Л}{{Л}}1 {М}{{М}}1 {Н}{{Н}}1 {О}{{О}}1 {П}{{П}}1 {Р}{{Р}}1
    {С}{{С}}1 {Т}{{Т}}1 {У}{{У}}1 {Ф}{{Ф}}1 {Х}{{Х}}1 {Ц}{{Ц}}1
    {Ч}{{Ч}}1 {Ш}{{Ш}}1 {Щ}{{Щ}}1 {Ъ}{{Ъ}}1 {Ы}{{Ы}}1 {Ь}{{Ь}}1
    {Э}{{Э}}1 {Ю}{{Ю}}1 {Я}{{Я}}1}
\lstset{style=labcode}
% Отдельный объект «код»: серый фон, поля 4 мм, перенос длинных строк
% и возможность продолжения на следующей странице без уменьшения шрифта.
\tcbset{lab code/.style={%
  listing only,listing engine=listings,breakable,
  % Иначе tcolorbox подавляет before skip сразу после заголовка.
  ignore nobreak=true,
  colback=labcodebg,colframe=labcodebg,boxrule=0pt,
  arc=0pt,outer arc=0pt,boxsep=0pt,
  left=\LabCodePadding,right=\LabCodePadding,
  top=\LabCodePadding,bottom=\LabCodePadding,
  before skip balanced=\LabCodeOuterSkip,
  after skip balanced=\LabCodeOuterSkip,
  left skip=10.7bp,right skip=3.8bp,
  width=\dimexpr\linewidth-14.5bp\relax,
  fontupper=\normalfont\normalsize}}
\newtcblisting{labcode}[1][]{lab code,listing options={style=labcode,#1}}
\newcommand\LabCodeFile[2][]{%
  \tcbinputlisting{lab code,listing file={#2},listing options={style=labcode,#1}}}
